%==============================================================================
\section{Học tăng cường sâu (Deep Reinforcement Learning)}
\label{sec:deep-reinforcement-learning}
%==============================================================================

\subsection{Học tăng cường sâu là gì?}

\begin{dinhnghia}[title={Deep Reinforcement Learning (Deep RL)}]
Học tăng cường sâu là nhánh con của học máy, kết hợp \textbf{học tăng cường} với \textbf{học sâu}.
Deep RL giải quyết bài toán cho phép agent tính toán học ra quyết định bằng cách tích hợp học sâu từ dữ liệu đầu vào \textbf{không có cấu trúc}, mà \textbf{không cần} kỹ thuật thủ công để thiết kế không gian trạng thái.
Các thuật toán Deep RL có khả năng quyết định hành động nào nên thực hiện để tối ưu target ngay cả khi đầu vào có kích thước lớn.
\end{dinhnghia}

\subsection{Các khái niệm then chốt}

\begin{ynghia}[title={Khối xây dựng Deep RL}]
Các khối xây dựng của học tăng cường sâu gồm mọi khía cạnh hỗ trợ học và agent ra quyết định.
Môi trường hiệu quả được tạo ra nhờ sự phối hợp của các yếu tố sau.
\end{ynghia}

\begin{dinhnghia}[title={Agent (Tác tử)}]
\textbf{Agent} là thực thể \textbf{học} và \textbf{ra quyết định}, tương tác với môi trường.
Agent hành động theo \textbf{chính sách} (policy) và tích lũy \textbf{kinh nghiệm}.
\end{dinhnghia}

\begin{dinhnghia}[title={Environment (Môi trường)}]
\textbf{Môi trường} là hệ thống bên ngoài agent mà agent giao tiếp.
Môi trường phản hồi cho agent dưới dạng \textbf{phần thưởng} hoặc \textbf{hình phạt} tùy theo hành động của agent.
\end{dinhnghia}

\begin{dinhnghia}[title={State (Trạng thái)}]
\textbf{Trạng thái} biểu diễn tình huống hoặc điều kiện \textbf{hiện tại} của môi trường tại một thời điểm, là cơ sở để agent đưa ra quyết định.
\end{dinhnghia}

\begin{dinhnghia}[title={Action (Hành động)}]
\textbf{Hành động} là lựa chọn agent thực hiện, làm \textbf{thay đổi trạng thái} của hệ thống.
\end{dinhnghia}

\begin{dinhnghia}[title={Policy (Chính sách)}]
\textbf{Chính sách} là kế hoạch định hướng ra quyết định của agent, \textbf{ánh xạ} trạng thái sang hành động.
\end{dinhnghia}

\begin{dinhnghia}[title={Value Function (Hàm giá trị)}]
\textbf{Hàm giá trị} ước lượng \textbf{phần thưởng tích lũy kỳ vọng} mà agent có thể đạt được từ một trạng thái cho trước khi theo một chính sách cụ thể.
\end{dinhnghia}

\begin{dinhnghia}[title={Model (Mô hình)}]
\textbf{Mô hình} biểu diễn \textbf{động học} của môi trường, cho phép agent mô phỏng kết quả tiềm năng của các cặp hành động--trạng thái phục vụ \textbf{lập kế hoạch}.
\end{dinhnghia}

\begin{phuongphap}[title={Chiến lược khám phá--khai thác (Exploration--Exploitation)}]
\textbf{Chiến lược khám phá--khai thác} là cách tiếp cận ra quyết định \textbf{cân bằng} giữa:
\begin{itemize}
  \item \textbf{Khám phá} hành động mới để học thêm;
  \item \textbf{Khai thác} hành động đã biết để nhận phần thưởng ngay lập tức.
\end{itemize}
\end{phuongphap}

\begin{dinhnghia}[title={Learning Algorithm (Thuật toán học)}]
\textbf{Thuật toán học} là phương pháp agent \textbf{cập nhật} hàm giá trị hoặc chính sách dựa trên kinh nghiệm thu được khi tương tác với môi trường.
\end{dinhnghia}

\begin{phuongphap}[title={Experience Replay (Bộ nhớ kinh nghiệm)}]
\textbf{Experience replay} là kỹ thuật \textbf{lấy mẫu ngẫu nhiên} từ các kinh nghiệm đã lưu trước đó trong quá trình huấn luyện, nhằm tăng \textbf{ổn định} học và giảm tương quan giữa các sự kiện liên tiếp.
\end{phuongphap}

\subsection{Cách hoạt động của Deep RL}

\begin{phuongphap}[title={Mạng nơ-ron và học thử--sai}]
Học tăng cường sâu dùng \textbf{mạng nơ-ron nhân tạo}, gồm các lớp nút mô phỏng hoạt động của nơ-ron trong não người.
Các nút xử lý và truyền thông tin qua phương pháp \textbf{thử--sai} để xác định kết quả hiệu quả.
\end{phuongphap}

\begin{ynghia}[title={Chính sách, tìm kiếm và đại diện}]
Trong Deep RL, thuật ngữ \textbf{policy} chỉ chiến lược máy tính hình thành dựa trên phản hồi nhận được khi tương tác môi trường.
Các policy này giúp máy tính ra quyết định bằng cách xét \textbf{trạng thái hiện tại} và \textbf{tập hành động} (các lựa chọn khả dụng).
Khi chọn các lựa chọn đó, diễn ra quá trình gọi là \textbf{search} (tìm kiếm): máy tính đánh giá hành động khác nhau và quan sát kết quả.
Khả năng phối hợp \textbf{học}, \textbf{ra quyết định} và \textbf{biểu diễn} có thể mang lại hiểu biết mới về cách não người vận hành.
\end{ynghia}

\begin{tinhchat}[title={Kiến trúc đặc trưng của Deep RL}]
Điểm tách biệt học tăng cường sâu nằm ở \textbf{kiến trúc}, cho phép học tương tự não người.
Kiến trúc gồm \textbf{nhiều lớp} mạng nơ-ron đủ mạnh để xử lý dữ liệu \textbf{không nhãn} và \textbf{không có cấu trúc}.
\end{tinhchat}

\subsection{Danh sách thuật toán Deep RL}

\begin{phuongphap}[title={Các thuật toán quan trọng}]
Dưới đây là một số thuật toán quan trọng trong học tăng cường sâu:
\begin{enumerate}
  \item \textbf{Deep Q-Network} hoặc \textbf{Deep Q-Learning}
  \item \textbf{Double Deep Q-Learning}
  \item \textbf{Actor--Critic Method}
  \item \textbf{Deep Deterministic Policy Gradient (DDPG)}
\end{enumerate}
\end{phuongphap}

\subsection{Ứng dụng của Deep RL}

\begin{ynghia}[title={Các lĩnh vực ứng dụng nổi bật}]
Học tăng cường sâu được dùng rộng rãi trong nhiều lĩnh vực sau.
\end{ynghia}

\subsubsection{Trò chơi (Gaming)}

\begin{vidu}[title={Trò chơi điện tử}]
Deep RL được dùng phát triển trò chơi vượt xa khả năng con người thông thường.
Các trò chơi thiết kế bằng Deep RL gồm game Atari 2600, cờ \textbf{Go}, \textbf{Poker} và nhiều thể loại khác.
\end{vidu}

\subsubsection{Điều khiển robot (Robot Control)}

\begin{phuongphap}[title={Học tăng cường đối kháng mạnh}]
Lĩnh vực này dùng \textbf{học tăng cường đối kháng mạnh} (robust adversarial reinforcement learning): agent học vận hành khi có đối thủ gây \textbf{nhiễu} lên hệ thống.
target là phát triển chiến lược tối ưu để xử lý gián đoạn.
\end{phuongphap}

\begin{vidu}[title={Robot trí tuệ nhân tạo}]
Robot do AI điều khiển có ứng dụng rộng: sản xuất, tự động hóa chuỗi cung ứng, y tế và nhiều lĩnh vực khác.
\end{vidu}

\subsubsection{Xe tự lái (Self-driving Cars)}

\begin{ynghia}[title={Lái tự động}]
Học tăng cường sâu là một trong các khái niệm then chốt của \textbf{lái tự động}.
Kịch bản lái tự động đòi hỏi hiểu môi trường, tác nhân tương tác, đàm phán và ra quyết định \textbf{động}---điều chỉ có được nhờ học tăng cường.
\end{ynghia}

\subsubsection{Y tế (Healthcare)}

\begin{vidu}[title={Chăm sóc sức khỏe cá nhân hóa}]
Deep RL đã thúc đẩy nhiều tiến bộ trong y tế, ví dụ \textbf{cá nhân hóa} liều thuốc để tối ưu chăm sóc bệnh nhân, đặc biệt với bệnh \textbf{mãn tính}.
\end{vidu}

\subsection{So sánh RL và Deep RL}

\begin{tinhchat}[title={Bảng so sánh Reinforcement Learning và Deep Reinforcement Learning}]
\begin{tabularx}{\linewidth}{>{\raggedright\arraybackslash}p{3.2cm}XX}
\textbf{Tiêu chí} & \textbf{Reinforcement Learning (RL)} & \textbf{Deep Reinforcement Learning (Deep RL)} \\
\midrule
Định nghĩa &
Nhánh con học máy dùng phương pháp \textbf{thử--sai} để ra quyết định. &
Nhánh con của RL tích hợp \textbf{học sâu} cho quyết định phức tạp hơn. \\
Xấp xỉ hàm (Function Approximation) &
Dùng phương pháp đơn giản như \textbf{bảng tra} (tabular) để ước lượng giá trị. &
Dùng \textbf{mạng nơ-ron} ước lượng giá trị, cho phép biểu diễn phức tạp hơn. \\
Biểu diễn trạng thái &
Phụ thuộc đặc trưng \textbf{thiết kế thủ công} để mô tả môi trường. &
\textbf{Tự học} đặc trưng liên quan từ dữ liệu đầu vào thô. \\
Độ phức tạp &
Hiệu quả với môi trường \textbf{đơn giản}, không gian trạng thái/hành động nhỏ. &
Hiệu quả với môi trường \textbf{chiều cao}, phức tạp. \\
Hiệu năng &
Tốt ở môi trường đơn giản; khó với không gian \textbf{lớn} hoặc \textbf{liên tục}. &
Vượt trội ở tác vụ phức tạp: game video, điều khiển robot. \\
Ứng dụng &
Tác vụ cơ bản như game đơn giản. &
Ứng dụng nâng cao: lái tự động, chơi game, điều khiển robot. \\
\end{tabularx}
\end{tinhchat}
